% Cato Maior de Senectute --- headnote
% cato-maior-de-senectute, 44 BC, Rome

\textit{\textit{Cato Maior de Senectute} --- Cato the Elder, On Old
Age --- was written in 44 BC, in the same extraordinary burst of
philosophical composition that followed Caesar's assassination and the
death of Cicero's daughter Tullia. Cicero was sixty-two; he dedicated
the book to his oldest friend, Titus Pomponius Atticus, then sixty-five,
and says plainly in the preface that he meant the consolation as much
for himself as for him --- a way for both of them to lighten the burden
of an old age now pressing or close at hand. He cast the argument as a
dialogue and put it not in his own mouth but in that of Marcus Porcius
Cato the Censor, the austere and famously vigorous old Roman, who
carries an authority no living speaker could. The dramatic date is 150
BC, the year before Cato's death at eighty-four; his listeners are the
two young men who would become the great figures of the next
generation, Scipio Aemilianus and Gaius Laelius.}

\textit{The structure is a clean refutation. Cato takes up the four
charges commonly brought against old age --- that it withdraws us from
active life, that it weakens the body, that it strips us of pleasure,
and that it brings us to the edge of death --- and answers each in turn.
Great affairs, he argues, are managed by counsel, authority, and
judgment, which age increases rather than diminishes; the loss of bodily
strength is met by keeping body and mind in use, and is more than repaid
by the mind's continued growth; the fading of appetite is no loss but a
release from a hard master, freeing a man for the higher and more
lasting pleasures, above all the satisfactions of farming, to which
Cato turns in the work's most famous and affectionate pages. The
nearness of death he meets last and most gravely: death is either
nothing to us or a passage to something better, the young die as surely
as the old, and a life rounded to its term falls due like ripe fruit.
The dialogue closes on the immortality of the soul --- Cato's
conviction, drawn from Plato and from the deathbed words of Cyrus in
Xenophon, that the soul is divine and undying, and his longing to
rejoin the friends who have gone before, the elder and younger
Scipios above all.}

\textit{Few of Cicero's works have been so steadily loved. Its blend of
Stoic and Peripatetic consolation, carried not by syllogism but by
Roman example and the warmth of a particular voice, made it a companion
to readers for two thousand years --- praised by Montaigne, set in type
early and often, and printed in America by Benjamin Franklin, who called
it a book for the comfort of the aged. Beneath the genial surface lies
the same private urgency that runs through all the works of these last
years: a man writing in danger and in grief, persuading himself, in the
voice of a sterner Roman than he was, that the final stretch of a
well-spent life is not a thing to be feared but the ripening and
crown of all the rest.}
